\section*{Suggested Improvements Based on Reviewers' Feedback}

Based on the reviewers' constructive comments, we propose the following improvements to strengthen the manuscript:

\subsection*{1. Discussion Section (Reviewer \#2, \#3)}
\begin{itemize}
    \item Add a dedicated \textbf{Discussion} section (after Results, before Conclusion).
    \item Address the \textbf{domain gap} between natural image/video
      pre-training (ImageNet, Kinetics-400) and thermal medical
      imaging. Discuss potential limitations and future directions
      such as domain adaptation, self-supervised pre-training on
      thermal data, or synthetic data augmentation.
    \item Critically analyze \textbf{overfitting risks} due to the
      small dataset size (149 patients). Discuss mitigation strategies
      such as ensembling, probability averaging, dropout, and data
      augmentation.
    \item Explicitly discuss the \textbf{clinical implications} of the
      results: whether dynamic thermal imaging is necessary, and
      whether large transformer models are justified compared to
      simpler architectures.
    \item Expand on \textbf{computational cost} of ViViT compared to
      ViT and 3D CNNs, and its implications for clinical deployment.
\end{itemize}

\subsection*{2. Dataset and Temporal Acquisition Details (Reviewer \#3)}
\begin{itemize}
\item Report the \textbf{frame rate (fps)}, total acquisition
  duration, and timing of cooling relative to capture in the DIT
  protocol.
    \item Clarify whether the 20 frames span the full physiological cooling response.
    \item Add ablation experiments:
    \begin{enumerate}
        \item Vary frame count (e.g., 3, 4, 6, 8, 20).
        \item Compare contiguous vs. uniformly spaced frames.
        \item Compare early vs. late frames.
    \end{enumerate}
  \item Discuss whether temporal modeling provides measurable
    improvements beyond spatial information.
\end{itemize}

\subsection*{3. Clarity and Presentation (Reviewer \#2, \#5)}
\begin{itemize}
\item Revise the \textbf{abstract} to minimize abbreviations; only
  retain those repeated multiple times (e.g., DIT, ViT).
    \item Add a \textbf{summary table} listing:
    \begin{itemize}
        \item Model type (ViT, ViViT, 3D CNN, etc.).
        \item Pre-training dataset (ImageNet-1K, Kinetics-400, or none).
        \item Fine-tuning dataset (DMR-IR, 149 patients).
        \item Input format (thermal maps vs. RGB).
    \end{itemize}
  \item Improve \textbf{Figure 1}: replace the flowchart with a more
    graphical schematic using standard deep learning symbols
    (convolutions, transformer blocks, temporal arrows).
  \item Ensure \textbf{references} are consistently ordered
    numerically within brackets (e.g., [17,21] not [21,17]).
  \item Verify all \textbf{ORCID IDs} of the authors.
\end{itemize}

\subsection*{4. Statistical Rigor (Reviewer \#5)}
\begin{itemize}
\item Add \textbf{confidence intervals} for AUC, F1, and accuracy
  metrics.
\item Perform \textbf{statistical significance tests} (e.g.,
  permutation tests, McNemar’s test) when comparing models.
\item Discuss suspiciously high cross-validation results (e.g., F1 =
  1.0 for ViT-32 pre-trained) and analyze whether this indicates
  overfitting.
\end{itemize}

\subsection*{5. Additional Discussion Points (Reviewer \#2, \#3)}
\begin{itemize}
\item Discuss the trade-off between \textbf{dynamic vs. static thermal
    imaging}: is the temporal signal essential, or are spatial
  features sufficient?
\item Compare \textbf{plain temperature matrices vs. infrared images}
  in terms of performance, acquisition complexity, and clinical
  practicality.
\item Highlight the need for \textbf{external validation} on larger or
  multicenter datasets to confirm generalizability.
\end{itemize}

\subsection*{6. Minor Revisions}
\begin{itemize}
    \item Justify the choice of sampling only 8 frames for TimeSformer (Reviewer \#5).
    \item Add limitations regarding dataset size, computational cost, and generalization.
    \item Ensure consistent terminology throughout the manuscript.
\end{itemize}

\section*{Summary}
In summary, the main revisions include: (i) adding a comprehensive
Discussion section, (ii) clarifying dataset acquisition details, (iii)
improving clarity with tables and figures, (iv) strengthening
statistical rigor, and (v) expanding on clinical implications and
limitations. These changes will enhance both the scientific and
clinical impact of the paper.

\section{Discussion}

This work presents, to our knowledge, the first application of Video
Vision Transformers (ViViTs) to dynamic infrared thermography (DIT)
for breast cancer classification. The proposed approach achieves
competitive performance (AUC $=0.9350 \pm 0.0094$ in cross-validation
and $0.9444$ on a held-out test set) and slightly advances prior DIT
baselines based on 3D CNNs. Below, we discuss the implications of
these findings, the necessity of dynamic acquisition and large-scale
transformer models, and the practical constraints that affect clinical
translation.

\subsection{Do we need dynamic thermal imaging?}
DIT aims to capture physiological responses under thermal stress
(e.g., vascular reactivity, metabolic heat signatures) that may not be
visible in single static frames. Our results show that a pre-trained
spatial ViT with temporal averaging can reach performance comparable
to a pre-trained video transformer, suggesting that a substantial
portion of the discriminative signal may be spatial. This raises two
complementary questions: (i) whether the current DIT timing and
sampling protocol sufficiently captures the temporal dynamics it is
designed to measure, and (ii) whether simpler static protocols (SIT)
augmented with robust spatial modeling could be adequate for certain
clinical scenarios.

To address this, future ablations should quantify the incremental
value of temporal modeling by varying (a) the number of frames (e.g.,
3, 4, 6, 8, and the full 20-frame sequence), (b) the sampling strategy
(contiguous vs.\ uniformly spaced frames), and (c) the timing (early
vs.\ late frames relative to cooling). Reporting frame rate, total
acquisition time, and stress onset/offset will allow the community to
relate model gains to physiological phases and decide when DIT
provides clinically meaningful information beyond spatial cues.

\subsection{Do we need large transformer models?}
Pre-training clearly benefited performance in our experiments (e.g.,
ViT-32 pre-trained vs.\ scratch), despite the domain gap between
natural images/videos and medical thermal data. However, large
transformers incur substantial computational cost, memory demand, and
latency, which can limit deployment in routine screening. The observed
parity between ViT-32 (spatial) and TimeSformer (spatio-temporal)
suggests a careful trade-off: in settings where temporal dynamics add
marginal gains, lighter spatial backbones or distilled variants may be
preferred.

Future work should evaluate: (i) compact architectures (e.g.,
mobile-friendly transformers or efficient CNN/Transformer hybrids),
(ii) knowledge distillation from ViViT into lightweight students, and
(iii) model pruning/quantization for point-of-care
devices. Benchmarking inference time, memory, and energy alongside
accuracy will enable more informed clinical choices.

\subsection{Domain gap and generalization}
Using ImageNet/Kinetics pre-training for thermal medical imaging
introduces a domain gap. While transfer learning improved results, it
may also encourage shortcut learning on spurious correlates. To reduce
this risk, several directions are promising: (i) self-supervised or
masked modeling pre-training directly on thermal sequences, (ii)
domain adaptation from natural to thermal imagery, and (iii) synthetic
data generation that respects thermal physics and acquisition
variability.

Generalization remains constrained by dataset size (149 patients for
IR; fewer for temperature maps). We recommend external validation on
larger and multi-center cohorts, reporting site-wise performance to
assess robustness to device, protocol, and population differences.

\subsection{Overfitting, statistical rigor, and evaluation design}
Near-perfect cross-validation outcomes (e.g., F1 $=1.0$ for ViT-32
pre-trained) warrant a cautious interpretation. Small patient-level
datasets and protocol homogeneity can inflate estimates. To strengthen
rigor, we suggest: (i) reporting confidence intervals for AUC, F1, and
accuracy; (ii) applying statistical tests such as permutation testing
or McNemar’s test when comparing paired model predictions; and (iii)
preferring patient-level ensembling or fold-wise probability averaging
over selecting a single best-performing fold/model.

Additionally, clarifying and standardizing the split strategy
(stratified, patient-level, non-overlapping sequences) helps avoid
leakage and improves reproducibility. Any data quality issues (e.g.,
resized frames, duplicated sequences) should be documented and
mitigated.

\subsection{Thermal maps vs.\ infrared images}
We observed that infrared RGB images delivered stronger performance
than raw temperature matrices. Possible explanations include: (i)
pre-trained backbones are tuned to image statistics closer to RGB
thermograms than to scalar temperature maps, and (ii) color encodings
may enhance contrast of physiologically relevant patterns for learned
filters. Nonetheless, temperature matrices provide physically grounded
measurements and may be preferable when interpretability and
calibration are paramount.

A practical path forward is to evaluate: (i) bespoke architectures for
temperature matrices with physics-informed inductive biases, and (ii)
cross-modal fusion (e.g., jointly modeling RGB thermograms and
temperature maps) to combine interpretability with discriminative
power.

\subsection{Clinical implications and deployment}
For clinical adoption, the value proposition must balance accuracy,
throughput, cost, and ease of integration. If temporal modeling adds
limited marginal gains, static protocols with efficient spatial models
could simplify workflows and reduce acquisition time. Conversely, if
ablations show that specific temporal windows or dynamics improve
sensitivity (e.g., in dense tissue or certain tumor phenotypes), DIT
with targeted sampling may be justified.

Interpretability is another key factor. Attention visualizations and
region-wise attribution over thermograms can help clinicians relate
model decisions to thermal asymmetries and vascular patterns. Finally,
clear reporting of computational requirements, training/inference
times, and hardware constraints will guide deployment decisions in
diverse clinical environments.

\subsection{Limitations and future work}
Our study is limited by: (i) small sample size; (ii) lack of explicit
timing metadata in DIT sequences; (iii) reliance on
natural-image/video pre-training; and (iv) absence of statistical
significance testing in the current comparisons. Future work should:
(a) include detailed temporal acquisition parameters and targeted
temporal ablations; (b) pursue domain-specific pre-training and
adaptation; (c) conduct external, multi-center validation; (d)
benchmark lightweight and distilled models; and (e) integrate
interpretability and uncertainty estimation to support clinical
decision-making.

Overall, while ViViTs show promise for DIT-based breast cancer
classification, the necessity of dynamic acquisition and large models
should be established through carefully controlled temporal studies,
rigorous statistical analyses, and deployment-oriented evaluations.
%%% Local Variables:
%%% mode: LaTeX
%%% TeX-master: "improvementIdeas.tex"
%%% End:
